Segmenting the lesion is the first step of most computer-aided dermoscopy pipelines: the mask defines the region from which border irregularity, asymmetry and colour variegation are measured, and it restricts classifiers to the lesion~\cite{mirikharaji2023survey}. Dermoscopic images, however, rarely show only skin and lesion. Terminal hairs cross the lesion border, rulers and calibration marks lie along the image edge, surgeons draw ink circles around suspicious lesions, and the circular aperture of the dermatoscope leaves dark corners. These structures are dark, thin or high-contrast, and they resemble lesion boundaries closely enough to pull a segmentation network away from the true contour. They also carry dataset-specific correlations: surgical skin markings were associated with altered melanoma predictions of a convolutional classifier~\cite{winkler2019markings}, and models trained on ISIC data can exploit artifacts as shortcuts~\cite{bissoto2020debiasing,geirhos2020shortcut}.

Three families of remedies exist. The oldest removes hair before analysis with morphological filtering and inpainting, as in DullRazor~\cite{lee1997dullrazor,telea2004inpainting}; such preprocessing is tuned by hand, handles one artifact type, and may erase genuine dark structures inside the lesion. Detection pipelines still apply it as a separate step~\cite{rana2025explainable}. A second family builds an artifact estimate into the network with fixed rules: FSSM-ResNet scores hair-like pixels by darkness, thinness and linearity in a fuzzy module and suppresses them during decoding~\cite{raikhankyzy2026fssm}. A third family learns the estimate end to end without artifact labels, either as a confidence map that scales a residual stream~\cite{singh2025dualresolution} or as multi-cue artifact responses that gate skip connections~\cite{gong2026artifactaware}. None of these methods tells the network which pixels are artifacts and of which type; the artifact map is either hand-crafted or left to emerge from the segmentation loss.

Supervising an artifact map directly has been impractical because pixel-level artifact annotations are scarce. Image-level artifact labels exist for the ISIC 2018 training set~\cite{bissoto2020debiasing,bissoto2022artifact}, a hair-mask set covers 500 images~\cite{hossain2023hairmask}, and marker and ruler masks exist for a small clinical set~\cite{carrion2023fedd}. Recent simulators show that artifacts can instead be generated with exact masks~\cite{kim2024ssynth,sizikova2026oasis}. We follow that route and combine it with partial labels: synthetic artifacts provide positives with exact masks, image-level labels and conservative colour and morphology rules provide negatives, and every remaining pixel is ignored by the loss.

This paper presents ArtiGate-YOLO, built on YOLO26n-seg~\cite{jocher2026yolo26} with an additional P2 detection level for small structures. Our earlier conference model added CBAM~\cite{woo2018cbam} to the stride-4 and stride-8 backbone features~\cite{iceeis2026own}; ArtiGate-YOLO keeps that network and changes only how attention is computed. The contributions are:
\begin{enumerate}
  \item An artifact head inside a real-time YOLO segmentation network that predicts four artifact types (hair, ruler, marker ink, dark corner) at stride~4 and is trained without manual artifact annotation, through procedurally composited artifacts with pixel-exact masks and a partial-label loss.
  \item Artifact-Gated Attention (AGA), a drop-in replacement for CBAM whose spatial attention receives the predicted artifact map and whose output is attenuated where artifacts are predicted. AGA reduces exactly to CBAM at initialization, and control rows separate the effects of training on artifact-augmented images, of the auxiliary artifact task and of gating with a supervised rather than an unsupervised map.
  \item A robustness evaluation by artifact type and severity: results stratified by artifact presence on the ISIC 2018 test set (image-level labels we release), a severity sweep with held-out synthetic artifacts, external validation on three test sets, and a test of the artifact maps against real marker, ruler and hair masks that were never used for training, all with image-level paired statistics over three training seeds.
\end{enumerate}

Section~\ref{sec:related} reviews related work, Section~\ref{sec:method} describes the model and its training targets, Section~\ref{sec:protocol} the datasets, baselines and statistical protocol, and Section~\ref{sec:results} the results. Section~\ref{sec:discussion} discusses limitations, and Section~\ref{sec:conclusion} concludes. The Supplementary Material holds the complete tables, the generator parameters and implementation details.
